\section{Basic Properties}

\begin{theorem}[Characterisation of Markov Chains]
    \label{thm:markov-chain-characterisation}%
    \(X \sim \Markov\pqty{\lambda, P}\) if and only if
    \[
        \Prob\pqty{X_0 = x_0, X_1 = x_1, \ldots, X_n = x_n} = \lambda_{x_0} P_{x_0 x_1} P_{x_1 x_2} \cdots P_{x_{n - 1} x_n}
    \]
    for all paths \(n \geq 0\) and \(x_0, x_1, \ldots, x_n \in I\), and
    \[
        \Prob\pqty{X_0 = i} = \lambda_i, \quad \forall i \in I.
    \]
\end{theorem}

\begin{proof}
    If the process \(X \sim \Markov\pqty{\lambda, P}\), then
    \begin{align*}
        \Prob\pqty{X_0 = x_0, \ldots, X_n = x_n} &= \Prob\pqty{\Conditional{X_0 = x_0}{X_{n - 1} = x_{n - 1}, \ldots, X_0 = x_0}} \cdot \Prob\pqty{X_{n - 1} = x_{n - 1}, \ldots, X_0 = x_0} \\
        &= P_{x_{n - 1} x_n} \Prob\pqty{X_{n - 1} = x_{n - 1}, \ldots, X_0 = x_0} \\
        &= \cdots \\
        &= P_{x_{n - 1} x_n} P_{x_{n - 2} x_{n - 1}} \cdots P_{x_0 x_1} \Prob\pqty{X_0 = x_0} \\
        &= P_{x_{n - 1} x_n} P_{x_{n - 2} x_{n - 1}} \cdots P_{x_0 x_1} \lambda_{x_0}.
    \end{align*}

    Conversely, if we have the above equation, then
    \begin{align*}
        \Prob\pqty{\Conditional{X_{n + 1} = x_{n + 1}}{X_n = x_n, \ldots, X_0 = x_0}} &= \frac{\Prob\pqty{X_{n + 1} = x_{n + 1}, X_n = x_n, \ldots, X_0 = x_0}}{\Prob\pqty{X_n = x_n, \ldots, X_0 = x_0}} \\
        &= \frac{\lambda_{x_0} P_{x_0 x_1} \cdots P_{x_{n - 1} x_n} P_{x_n x_{n + 1}}}{\lambda_{x_0} P_{x_0 x_1} \cdots P_{x_{n - 1} x_n}} \\
        &= P_{x_n x_{n + 1}} \\
        &= \Prob\pqty{\Conditional{X_{n + 1} = x_{n + 1}}{X_n = x_n}}.
    \end{align*}
\end{proof}

\begin{remark}
    If \(X = \pqty{X_n} _{n \geq 0}\) is a Markov chain, then given \(X_{n}\), \(X_{n + 1}\) and \(X_{n - 1}\) are \emph{conditionally independent}, but not typically independent.
\end{remark}

\begin{example}
    If \(\pqty{X_n}_{n \geq 0}\) are independent and identically distributed, then \(\pqty{X_n}\) is a Markov chain with transition matrix \(P_{xy} = \Prob\pqty{X_1 = y}\) for all \(x, y \in I\).
\end{example}

\begin{definition}[Unit Mass Notation]
    For any \(i \in I\), we let \(\kronecker_i\) denote the probability distribution that places a unit mass at \(i\):
    \[
        \delta_{ij} = \indicator\Bqty{j = i} = \begin{cases}
            1, & j = i, \\
            0, & j \neq i.
        \end{cases}
    \]
\end{definition}

Recall the following facts.

\begin{itemize}
    \item The random variables \(\pqty{X_0, X_1, \ldots, X_n}\) are \emph{independent}, if
    \[
        \Prob\pqty{X_0 = x_0, X_1 = x_1, \ldots, X_n = x_n} = \Prob\pqty{X_0 = x_0} \Prob\pqty{X_1 = x_1} \cdots \Prob\pqty{X_n = x_n}.
    \]
    \item The random variables \(\pqty{X_0, X_1, \ldots}\) are \emph{independent}, if \(\pqty{X_0, X_1, \ldots, X_n}\) are independent, for all \(n \geq 0\).
    \item The processes \(X = \pqty{X_n}_{n \geq 0}\) and \(Y = \pqty{Y_n}_{n \geq 0}\) are \emph{independent}, if for any \(0 \leq t_1 < t_2 < \cdots < t_m\) and \(0 \leq s_1 < s_2 < \cdots < s_k\), we have
    \begin{align*}
        &\Prob\pqty{X_{t_1} = x_1, \ldots, X_{t_m} = x_m, Y_{s_1} = y_1, \ldots, Y_{s_k} = y_k} \\
        &\hspace{16em} = \Prob\pqty{X_{t_1} = x_1, \ldots, X_{t_m} = x_m} \Prob\pqty{Y_{s_1} = y_1, \ldots, Y_{s_k} = y_k}.
    \end{align*}
    for all \(x_1, \ldots, x_m, y_1, \ldots, y_k\).
\end{itemize}

\begin{theorem}[Markov Property]
    \label{thm:markov-property}%
    Suppose \(X = \pqty{X_n}_{n \geq 0} \sim \Markov\pqty{\lambda, P}\) and let \(m \geq 1\) and \(i \in I\). Then, conditional on \(\Bqty{X_m = i}\), the new process \(\pqty{X_{m + n}}_{n \geq 0} \sim \Markov\pqty{\delta_i, P}\) and it is independent of \(\pqty{X_0, \ldots, X_m}\)
\end{theorem}

\begin{proof}
    Clearly,
    \[
        \Prob\pqty{X_m = j} = \delta_{ij}
    \]
    since we conditioned on \(\Bqty{X_m = i}\).
    \begin{align*}
        &\phantom{=} \Prob\pqty{\Conditional{X_m = x_m, X_{m + 1} = x_{m + 1}, \ldots, X_{m + n} = x_{m + n}}{X_m = i}} \\
        &= \indicator \Bqty{x_m = i} \frac{\Prob\pqty{X_m = x_m, X_{m + 1} = x_{m + 1}, \ldots, X_{m + n} = x_{m + n}}}{\Prob\pqty{X_m = x_m}} \\
        &= \frac{\indicator \Bqty{x_m = i}}{\Prob\pqty{X_m = x_m}} \sum_{x_0, \ldots, x_{m - 1} \in I} \Prob\pqty{X_0 = x_0, \ldots, X_m = x_m, X_{m + 1} = x_{m + 1}, \ldots, X_{m + n} = x_{m + n}} \\
        &= \frac{\indicator \Bqty{x_m = i}}{\Prob\pqty{X_m = x_m}} \sum_{x_0, \ldots, x_{m - 1} \in I} \lambda_{x_0} P_{x_0 x_1} \cdots P_{x_{m - 1} x_m} P_{x_m x_{m + 1}} \cdots P_{x_{m + n - 1} x_{m + n}} \\
        &= \frac{\indicator \Bqty{x_m = i}}{\Prob\pqty{X_m = x_m}} P_{x_m x_{m + 1}} \cdots P_{x_{m + n - 1} x_{m + n}} \sum_{x_0, \ldots, x_{m - 1} \in I} \lambda_{x_0} P_{x_0 x_1} \cdots P_{x_{m - 1} x_m} \\
        &= \frac{\indicator \Bqty{x_m = i} \Prob\pqty{X_m = x_m}}{\Prob\pqty{X_m = x_m}} P_{x_m x_{m + 1}} \cdots P_{x_{m + n - 1} x_{m + n}} \\
        &= \indicator \Bqty{x_m = i} P_{x_m x_{m + 1}} \cdots P_{x_{m + n - 1} x_{m + n}}
    \end{align*}
    which is precisely the condition for \(\Markov\pqty{\delta_i, P}\), by \autoref{thm:markov-chain-characterisation}.

    To prove independence, fix times
    \[
        m < i_1 < i_2 < \cdots < i_k
    \]
    and states
    \[
        x_0, x_1, \ldots, x_m, x_{m + 1}, \ldots, x_{m + k}.
    \]

    Then,
    \begin{align*}
        \phantom{=} &\Prob\pqty{\Conditional{X_0 = x_0, \ldots, X_m = x_m, X_{i_1} = x_{m + 1}, \ldots, X_{i_k} = x_{m + k}}{X_m = i}} \\
        &= \indicator\Bqty{x_m = i} \frac{\Prob\pqty{X_0 = x_0, \ldots, X_m = x_m, X_{i_1} = x_{m + 1}, \ldots, X_{i_k} = x_{m + k}}}{\Prob\pqty{X_m = x_m}} \\
        &= \frac{\indicator\Bqty{x_m = i}}{\Prob\pqty{X_m = x_m}} \Prob\pqty{X_0 = x_0, \ldots, X_m = x_m} \Prob\pqty{\Conditional{X_{i_1} = x_{m + 1}, \ldots, X_{i_k} = x_{m + k}}{X_0 = x_0, \ldots, X_m = x_m}} \\
        &= \frac{\indicator\Bqty{x_m = i}}{\Prob\pqty{X_m = x_m}} \lambda_{x_0} P_{x_0 x_1} \cdots P_{x_{m - 1} x_m} \Prob\pqty{\Conditional{X_{i_1} = x_{m + 1}, \ldots, X_{i_k} = x_{m + k}}{X_m = x_m}}\\
        &= \Prob\pqty{\Conditional{X_0 = x_0, \ldots, X_m = x_m}{X_m = i}} \cdot \Prob\pqty{\Conditional{X_{i_1} = x_{m + 1}, \ldots, X_{i_k} = x_{m + k}}{X_m = i}}
    \end{align*}
    using \autoref{thm:markov-chain-characterisation} again, note that the definition of a Markov chain is applied in the third equal sign, more rigorously by summing through the possibilities of the intermediate times between \(i_1, \ldots, i_k\). This finishes the proof.
\end{proof}